\chapter{Getting Paid: EOR, Contracting, W-8BEN \& GST}
\label{ch:getting-paid}

Every chapter so far is about getting the offer. This short chapter covers what the source material flags right
after that point: how you actually get paid, legally, once a US or European company says yes to hiring you from
India. None of this is legal or tax advice — treat it as a map of the questions to ask a professional before you
sign anything.

\section{Two Ways a Foreign Company Can Hire You}

\begin{longtable}{p{0.28\textwidth} p{0.64\textwidth}}
\toprule
\textbf{Arrangement} & \textbf{What it means} \\
\midrule
\endhead
Employer of Record (EOR) & A third-party service (cited examples: Deel, Remote, Rippling) legally employs you on
behalf of the hiring company in your country, handling local compliance, payroll, and often benefits. You get an
employment relationship closer to a normal job, with the hiring company directing your work. \\
Independent contractor & You invoice the company directly for your services, with no employer--employee
relationship. You're responsible for your own tax filing, and the company is not. This is typically the simpler
arrangement for a company to set up, but shifts compliance work onto you. \\
\bottomrule
\end{longtable}

\begin{keytakeaway}
Which one you end up with isn't usually your choice to make unilaterally — it depends on what the hiring company
is set up to support. What \emph{is} your responsibility is understanding which one you're in, since it changes
your tax filing, your protections, and your paperwork.
\end{keytakeaway}

\section{The W-8BEN Form}

If you're hired as an independent contractor by a \textbf{US} company, you'll likely be asked to complete a
\textbf{Form W-8BEN}. Its purpose: it certifies to the IRS that you are a foreign person, not a US taxpayer, which
is what allows the US company to pay you without default US tax withholding on US-sourced income.

\begin{caution}
Whether withholding applies, and at what rate, depends on the type of income, whether a US--India tax treaty
provision applies, and how the engagement is structured — it is \emph{not} automatic that filing a W-8BEN means
zero US tax exposure. Don't assume the form alone settles the question; have a professional confirm your specific
situation, especially if the engagement involves anything beyond straightforward services income.
\end{caution}

\section{GST on Exported Services from India}

Freelance or contract income earned by an India-based professional for a foreign client is generally treated as
an \textbf{export of services} under India's GST framework, which can qualify for zero-rated GST treatment
(meaning no GST is charged on the invoice) — but this status is conditional, not automatic.

\begin{caution}
Zero-rating for exported services typically depends on meeting specific conditions: payment received in
convertible foreign exchange, correct invoicing, and (depending on turnover and structure) GST registration or a
Letter of Undertaking (LUT). The specifics change with policy updates and depend on your turnover and business
structure — verify current requirements with a chartered accountant rather than relying on a general summary,
including this one.
\end{caution}

\section{A Pre-Signing Checklist}

Before accepting a remote offer from a foreign company, get clear, ideally written, answers to:

\begin{itemize}
  \item Am I being hired as an employee (via an EOR) or as an independent contractor?
  \item If a contractor: will I need to file a W-8BEN (US) or an equivalent form for another country?
  \item How will I invoice, and in what currency?
  \item Do I need GST registration or an LUT for this engagement, given my current turnover and structure?
  \item Who is responsible for my health coverage, paid leave, and other benefits under this arrangement, if
  anyone?
  \item What does the contract say about IP ownership, termination notice, and payment terms (net-15, net-30,
  etc.)?
\end{itemize}

\begin{keytakeaway}
None of this should block you from pursuing the opportunities in the rest of this guide — it's a short checklist
to run \emph{once an offer is real}, not a reason to hesitate before that point. Budget for a short consultation
with a chartered accountant familiar with export-of-services income before your first contract payment arrives;
it's a small cost relative to getting the structure wrong on a recurring income stream.
\end{keytakeaway}

